\section*{はじめに}
\subsection*{授業のテーマ}

「見えないもの」をみよう・測ろう・考えようとするのが心理統計の世界である。

\subsection*{一年を通じて伝えたいポイント}
\begin{description}
	\item[母数を知るための推測] 大学での統計，心理統計は標本の記述統計量を検証したいのではなく，母集団を表す数字を求めるための推測統計である。
	\item[確率を用いた推論] 母数を知るための方法は3つある。代表値を用いたモーメント法，確率分布をつかった最尤法，ベイズ法である。
	\item[要因計画と線形モデル] 心理学では，要因計画による研究条件をコントロールし，平均値の差をもちいた考察(平均因果効果)をおこなうことが多い。そしてそれらは総じて線形モデルとして表現できる。
	\item[モデル比較と意思決定] モデルによる母数の推定だけではなく，そこから一定の「結論」あるいは「意思決定」を行うための方法として，モデル比較やNHSTといった方法がある。
	\item[統計環境Rによる実践] 統計環境Rを利用して，理論的な理解だけでなく実際に計算ができることも身につけるべき技術である。
\end{description}


なお，各コマにおける「標準シラバスにおける位置づけ」とは，公認心理師大学カリキュラム　標準シラバス(2018年8月22日版)\footnote{\url{https://psych.or.jp/wp-content/uploads/2018/04/standard_syllabus_2018-8-22.pdf}} との対応を表しています。

\subsection*{その他}
授業シラバスとこの講義資料を掲載したサイト(\url{https://kosugitti.github.io/psychometrics_syllabus/})で，最新版のシラバスと授業資料，授業で用いるサンプルデータやコードの配布を行なっています。
